\chapter*{Laboratory and Workshop Safety}
\addcontentsline{toc}{chapter}{Laboratory and Workshop Safety}

This chapter establishes mandatory safety practices for all EDNS~491/492 lab, shop, and project work. Capstone projects span electrical, mechanical, chemical, and computational domains. The hazards you encounter depend on your project, but the practices below apply universally.

\section*{EHS Training Requirement}
Both SDI and SDII require completion of an EHS Safety Training module on Canvas. In SDII, you must also complete a \textbf{Hazard Assessment and Safety Plan} specific to your project by Week~2. This document identifies the hazards associated with your build, the controls you will implement, and the PPE required. Your PA must review and approve the safety plan before fabrication begins. See Chapter~21 for the full EHS framework.

\section*{Electrical Hazards}
Electrical shock severity depends on current through the body, not voltage alone, but higher voltage drives more current through the body's resistance ($\sim$1\,k$\Omega$ hand-to-hand, wet skin). Below 30\,VDC or 25\,VAC RMS is generally safe for dry skin contact (IEC~61010). Between 30\,V and 48\,V, use insulated tools and the one-hand rule. Above 48\,V, including all AC mains: de-energize and lock out before touching. \textbf{Never design your own AC-DC converter.} Always use a pre-certified (UL/CE listed) wall adapter.

\section*{Machine Shop Safety}
Before using any machine in the Labriola InnoHub shops (Metal Shop, Wood Shop, Composites Shop), you must complete the machine-specific training on the Labriola InnoHub Canvas site. No drop-ins are accepted. CNC training requires prior completion of the corresponding manual machine training. Extended machine use ($>$4 hours) for mills, lathes, welders, wood lathe, and table saw requires manager approval. Training prioritization: \textbf{safety training always takes precedence over project work}. If a TA needs a machine or work area for training, you will be asked to return at a different time.

\section*{3D Printing Safety}
FDM printers produce fumes, especially with ABS and specialty filaments. Work in ventilated areas. SLA resin is a skin sensitizer and eye irritant: wear nitrile gloves and safety glasses when handling uncured parts. UV exposure from curing stations can cause eye damage; do not look directly at UV LEDs. Review the InnoHub 3D printing policies in Chapter~12 for print limits, material restrictions, and overnight print approvals.

\section*{Lithium Battery Safety}
Lithium-ion and lithium-polymer cells present fire and explosion risks if mishandled. Never short-circuit, overcharge, over-discharge, puncture, crush, or overheat cells. Always use a Battery Management System (BMS) for charging and discharge protection. If a cell swells, hisses, or smokes: clear the area, place in a fire-safe container (metal bucket with sand), and notify lab staff immediately. Never charge lithium cells unattended. Always charge in a fireproof LiPo safety bag.

\section*{Soldering Safety}
Soldering irons operate at 300--400\,\degC{}. Return the iron to its stand when not actively soldering. Use a fume extractor or work in a ventilated area. Wash hands after soldering, especially if using leaded solder.

\section*{Chemical and Coatings Safety}
Adhesives, solvents, paints, and coatings may be flammable, toxic, or sensitizing. Check the Safety Data Sheet (SDS) before use. Spray painting must use the \textbf{designated paint booth} in the InnoHub (complete the paint booth pre-quiz on Canvas first). Never spray paint in open lab areas.

\section*{Compressed Gas and Pneumatic Systems}
Compressed gas cylinders must be secured upright and chained to a wall or cart. Never direct compressed air at a person. Pneumatic systems store energy; relieve all pressure before disconnecting fittings.

\section*{Personal Protective Equipment (PPE)}
Minimum PPE varies by activity: safety glasses for all shop work, hearing protection near loud machinery, closed-toe shoes in all shop spaces, gloves for chemical handling (not for rotating machinery). Long hair must be tied back; remove dangling jewelry and loose clothing before operating any rotating equipment.

\section*{Emergency Procedures}
Know the location of the nearest fire extinguisher, first aid kit, eyewash station, and emergency exit \emph{before} beginning any work. In an emergency, call 911. For non-emergency EHS concerns, contact the InnoHub Operations Manager or Capstone Design TAs.


% ═══════════════════════════════════════════════════════════════
